Program 12.                   \documentclass[12pt]{report}
\usepackage{graphicx}
\usepackage{subcaption}

\begin{document}
	
%------------------ TITLE PAGE ------------------

\begin{titlepage}
\centering

\vspace*{0.5cm}

\LARGE{\bf VISVESVARAYA TECHNOLOGICAL UNIVERSITY}

\vspace{1cm}

\large{\textbf{GOVERNMENT ENGINEERING COLLEGE\\[0.2cm]
HASSAN – 573201\\[0.3cm]}}

\vspace{0.5cm}

\textbf{DEPARTMENT OF COMPUTER SCIENCE \& ENGINEERING}\\[5cm]

\textbf{\underline{\large Submitted By}}\\[0.3cm]

\begin{tabular}{l c l}

\textbf{Sindhu Raj} & \textbf{:} & \textbf{4GH21CS047}\\

\end{tabular}

\end{titlepage}

	
	\tableofcontents
	\listoffigures
	
	%------------------ CHAPTER 1 ------------------

\chapter{ Free Software}

\section{What is Free Software?}

The word free software refers to software that respects
users' freedom and community. Roughly, it means that
the users have the freedom to run, copy, distribute,
study, change and improve the software.

The term free software is sometimes misunderstood.
The word free in free software refers to freedom,
not price.

\section{The Free Software Definition}

The four essential freedoms are:

\begin{itemize}
\item The freedom to run the program as you wish,
for any purpose. (Freedom 0)

\item The freedom to study how the program works
and change it so it does your computing as you wish.
(Freedom 1)

\item The freedom to redistribute copies so you can
help others. (Freedom 2)

\item The freedom to distribute copies of your
modified versions to others. (Freedom 3)
\end{itemize}

By doing this, you can give the whole community
a chance to benefit from your changes.

%------------------ CHAPTER 2 ------------------

\chapter{Listing Environment}

\section{Unordered Lists}

\subsection{Groceries List}

\begin{itemize}
\item Eggs
\item Milk
\item Biscuits
\item Rice
\end{itemize}

\subsection{Football Teams}

\begin{itemize}

\item English Premier League

\begin{itemize}
\item Manchester United
\item Liverpool
\end{itemize}

\item La Liga

\begin{itemize}
\item Barcelona
\item Real Madrid
\end{itemize}

\item Bundesliga

\begin{itemize}
\item Bayern Munich
\item Borussia Dortmund
\end{itemize}

\end{itemize}

\section{Ordered Lists}

\subsection{ICC WTC Rankings}

\begin{enumerate}
\item India
\item Australia
\item New Zealand
\end{enumerate}

\subsection{Countries Ranked by Market Cap}

\begin{enumerate}

\item Asia

\begin{enumerate}
\item China
\item Japan
\item India
\end{enumerate}

\item Europe

\begin{enumerate}
\item United Kingdom
\item France
\item Germany
\end{enumerate}

\end{enumerate}
	
	\chapter{Methods}
	
	\section[Method's section]{This is Method's section}
		\subsection*{This is Method's subsection}
	
	\chapter{Results}
	\addcontentsline{toc}{chapter}{Result's section}
	This is the results section.
	
	\chapter*{snapshots}
	\addcontentsline{toc}{chapter}{Snaps}
	This is the Snapshot section.
	
	\begin{figure}[h]
	
		\begin{subfigure}{0.40\linewidth}
		\centering
		\includegraphics[width=\linewidth]{TeXworks.png}
		\caption{TexWorks}
		\end{subfigure}
		\hfill
		\begin{subfigure}{0.40\linewidth}
		\centering
		\includegraphics[width=\linewidth]{vtu_logo.png}
		\caption{Latex1}
		\end{subfigure}
		\caption{Latex and TeXWorks}
		
	\end{figure}
	
	\begin{figure}[h]
	\centering
	\includegraphics[width=4cm]{latex.png}
	\caption{Latex}
	\end{figure}
	
	\chapter{Conclusion}
	This is the conclusion.
\end{document}